\section{Introduction}

\subsection{Problem Statement and Motivation}

Artificial intelligence (AI) and, in particular, deep learning have enabled substantial advances in
medical signal analysis, supporting early detection, and continuous monitoring.

Epilepsy, a common chronic neurological disorder affecting roughly 7.6 per 1,000 people
\parencite{beghiEpidemiologyEpilepsy2019}, remains clinically challenging: about 30\% of
patients are drug resistant \parencite{kwanEarlyIdentificationRefractory2000,chenNewEraPersonalised2020},
while effective pharmacotherapy can entail adverse side effects that reduce quality of life.

New improvements in machine learning particularly in deep learning architectures open
up the possibility of major advancements in seizure detection and prediction for Epilepsy patients.

Generally, seizure detection and prediction offer the potential to enable timely
intervention (e.g., rescue medication), injury prevention, and improved safety, including
mitigation of risks associated with SUDEP and seizure‑related falls \parencite{hesdorfferCombinedAnalysisRisk2011}.

In hospital and research settings, high‑quality electroencephalography (EEG) enables
reliable detection/prediction, but EEG hardware (multiple electrodes, gel, maintenance)
is obtrusive, limiting pervasive ambulatory monitoring \parencite{chenNewEraPersonalised2020}.

Wearable, non‑EEG sensing (ECG/HRV, PPG, EDA, accelerometry) combined with modern
machine learning in particular deep learning architectures offers a path toward accessible, 
continuous seizure monitoring outside the clinic.

However, achieving clinically meaningful performance in real‑world environments faces key
constraints:

\begin{itemize}
    \item Requirement for high sensitivity without an impractical false alarm rate (precision/ specificity trade‑offs)

    \item Robustness to movement, physiological variability, and non‑stationarity

    \item Interpretability and safety considerations in regulated healthcare

    \item Differences between clinical (controlled) vs everyday (noisy, resource‑limited)
    deployment contexts

    \item Limited availability of prospective, patient‑preserving evaluation and multimodal
    wearable datasets
\end{itemize}

\subsection{Research Aim and Questions}

This literature review focuses on the use of non‑EEG, wearable/autonomic biosignals for epileptic
seizure detection and prediction and analyzes how model classes and study designs
influence performance across clinical and ambulatory contexts.

\vspace{1em}

\subsubsection{Primary Research Question (RQ):}

Which machine learning and deep learning model architectures applied to wearable autonomic
biosignals (ECG/HRV, PPG, EDA, ACC) achieve viable seizure detection and prediction performance (high sensitivity with acceptable false alarm rates)
under different study phases (retrospective, prospective) and deployment contexts
(clinic vs everyday use)?

\subsubsection{Sub‑questions}

\begin{enumerate}
    \item How do classical feature‑based models compare to newer deep multimodal architectures
    (e.g., 1D‑CNN, LSTM, attention/fusion) in robustness and generalization?

    \item What is the impact of study design (retrospective, pseudo‑prospective, prospective)
    and evaluation granularity (sample‑ vs alarm‑based) on reported performance metrics?

    \item How does patient‑specific fine‑tuning of the model affect sensitivity and false
    alarm rate in ambulatory scenarios?

    \item Which gaps (data availability, focal seizure coverage, prospective validation)
    limit translation, and how might augmentation or synthetic data help without
    introducing bias?
\end{enumerate}

\subsection{Methodological Approach and Structure}

The review will follow established guidance for structured IS/medical literature analyses
(concept matrix inspired by \cite{websterAnalyzingPastDetermine2002} and PRISMA principles
for transparency).

A multi‑stage search (2015--2025) will combine:

(a) scoping via Google Scholar;

(b) formal queries in IEEE Xplore, PubMed, and Scopus using controlled keyword strings;

and (c) backward/forward citation chasing on seminal papers.

Inclusion focuses on human studies using wearable/non‑EEG autonomic signals for seizure
detection or preictal prediction.

Exclusion criteria remove purely EEG‑only, invasive, or purely algorithmic
simulation studies.

Extracted dimensions will populate a concept matrix:

(1) setting (clinical vs ambulatory);

(2) modality (ECG/HRV, PPG, EDA, ACC);

(3) task (detection vs prediction);

(4) model class (classical feature‑based vs deep/multimodal fusion);

(5) study phase (retrospective, pseudo‑prospective, prospective);

(6) evaluation granularity (sample‑ vs alarm‑based);

(7) metrics (sensitivity, FAR/h, AUC, latency);

(8) personalization (patient‑specific adaptation);

(9) interpretability/safety considerations;

(10) data characteristics/ utilized dataset (cohort size, seizure count, class imbalance).

Quality flags include presence of patient‑preserving splits, external/cross‑dataset
validation, and reporting of false alarms per hour.

A PRISMA‑style flow diagram will summarize identification, screening, eligibility, and
inclusion.

Synthesis will be primarily narrative, grouped by modality and task, highlighting deployment
feasibility and gaps impeding translation.

\subsubsection{Preliminary Outline (approx. target pages for 15 pages of main text):}

\begin{enumerate}
    \item Introduction (Problem, Aim, Methodology) -- 2 pp

    \item Technical Background (Epilepsy, ANS, Wearables, ML basics) -- 2 pp

    \item Methodology (Search, Selection, Extraction, PRISMA flowgraph) -- 2 pp

    \item Results (Concept Matrix, Comparative Performance, Modality/Model
    Analysis) -- 5 pp

    \item Discussion (Interpretation, Trade‑offs, Safety, Deployment) -- 2 pp

    \item Limitations and Future Work -- 1 pp

    \item Conclusion -- 1 pp
\end{enumerate}

\section{Technical Background}

\subsection{Epileptic Seizures}
Epileptic seizures are classified by their onset as focal (starting in a specific area of the brain)
or generalized (affecting both hemispheres from the onset).
\begin{enumerate}
    \item \textbf{Focal Seizures}: Seizures are further split up into focal aware and focal impaired awareness seizures,
    depending on whether the patient is aware during the seizure or not, and by onset features (motor vs non‑motor).
    Focal seizures may evolve to bilateral tonic–clonic seizures.
    \item \textbf{Generalized Seizures}: Generalized seizures begin in both cerebral hemispheres and include motor
     (e.g. tonic–clonic, myoclonic, atonic) and non‑motor (absence) types.
\end{enumerate}
\parencite{fisherInstructionManualILAE2017}.
\textbf{Note:} Tonic‑clonic seizures start with a brief tonic phase of sustained muscle contraction followed
by a clonic phase of rhythmic jerking; they commonly cause loss of consciousness and carry increased risk of injury and SUDEP \parencite{hesdorfferCombinedAnalysisRisk2011}.

\subsection{Autonomic Nervous System}
The autonomic nervous system (ANS) comprises the sympathetic (arousal)
and parasympathetic/vagal (rest) branches. Heart‑rate variability (HRV)
reflects how these two balance each other: HF power and RMSSD mainly
indicate parasympathetic (calming) activity. LF includes a mix of both
branches. The LF/HF ratio is only a rough summary of this balance.
Electrodermal activity (EDA) is a direct sympathetic output with a tonic
skin conductance level (SCL) and phasic skin conductance responses (SCRs).
Photoplethysmography (PPG) yields pulse‑rate variability analogous to HRV
from ECG \parencite{mironAutonomicBiosignalsSeizure2025,masonHeartRateVariability2024}.

\subsection{Non-tonic–clonic seizures and autonomic markers}
Non‑convulsive/focal seizures typically elicit subtle autonomic signatures
rather than large motor patterns. Sympathetic activation manifests as:
\begin{enumerate}
    \item ictal (during a seizure) tachycardia and preictal HR rise with reduced
    HRV (lower RMSSD/HF, higher LF/HF) measurable from ECG or PPG‑derived
    pulse rate variability;
    \item increased electrodermal activity (tonic SCL, phasic SCRs);
    \item limited accelerometry changes mainly useful for artifact rejection and
    context (e.g., sleep or daily activity).
\end{enumerate}
Patient‑specific thresholds or fine‑tuning are often needed due to differences in autonomic
reactivity and medication effects between patients \parencite{sethFeasibilityCardiacbasedSeizure2023,masonHeartRateVariability2024}.

\subsection{Wearable Technology}
Wearable heart‑rate sensing shows high agreement with chest‑strap ECG under real‑world motion: Apple Watch Ultra achieves overall concordance of 0.998 with a Polar H10 reference, Garmin Enduro 2 0.995, while Polar Pacer Pro is lower at 0.957.
Accuracy degrades on the dominant wrist due to motion artifacts.
Chest‑strap ECG provides robust RR‑interval quality during exercise and is recommended as a reference when Holter signal quality degrades with activity \parencite{wolfHearttoWearAssessingAccuracy2024a}.

\subsection{Machine Learning}

Machine learning on wearable biosignals (ECG/HRV, EDA, ACC, PPG) typically frames seizure detection
and prediction as standard classification. Classical models (e.g., SVM, RF, kNN) on HRV/EDA/ACC features
achieve solid detection and some prediction, but often in retrospective designs
\parencite{sethFeasibilityCardiacbasedSeizure2023,masonHeartRateVariability2024}.

Deep architectures (1D‑CNN, LSTM) and multimodal fusion remain the most promising, improving robustness and phase classification,
especially when combining ECG/EDA/ACC with attention/fusion layers
\parencite{yangMultimodalAISystem2022,pordoyEnhancedNonEEGMultimodal2025}.

ECG‑focused studies add interpretable features and out‑of‑distribution evaluation but emphasize non‑stationarity,
class imbalance, and scarce prospective validation as key limitations
\parencite{abtahiIdentificationRelevantECG2025,kalousiosECGbasedEpilepticSeizure2024}.

\subsection{Patient-Specific Fine-Tuning}
Inter‑patient variability in autonomic responses and preictal dynamics is substantial;
the most informative features often differ by patient, which degrades population‑trained models.
Patient‑specific fine‑tuning typically improves sensitivity and lowers false alarms in ambulatory, alarm‑based settings
\parencite{andradePerformanceSeizurePrediction2024}.

\subsection{Datasets and Synthetic Data}
Public datasets for seizure detection/prediction are largely EEG‑centric with auxiliary ECG
(TUH, EPILEPSIAE, RPAH, CHB‑MIT), while truly wearable multimodal corpora are scarce;
one of the few open non‑EEG sources is the Open Seizure Database (ACC+HR)
\parencite{yangMultimodalAISystem2022,pordoyEnhancedNonEEGMultimodal2025}.
Cross‑dataset and pseudo‑prospective evaluation across TUH→RPAH/EPILEPSIAE and across EPILEPSIAE/CHB‑MIT/AES/Epilepsy
Ecosystem exposes domain shift and favors alarm‑based over sample‑based metrics
\parencite{yangMultimodalAISystem2022,andradePerformanceSeizurePrediction2024}.
Given rare events, imbalance, and noisy labels, studies use augmentation (time‑warping, jitter, scaling)
and synthetic data to balance preictal/ictal segments, with patient‑preserving splits and external validation to
avoid device/site overfitting
\parencite{kalousiosECGbasedEpilepticSeizure2024,sethFeasibilityCardiacbasedSeizure2023}.

\section{Methodology}
\begin{enumerate}
    \item \textbf{Scope:} human studies in journals (2015--2025) on non‑EEG wearable/autonomic biosignals (ECG/HRV, PPG, EDA, ACC) for seizure detection or prediction.
    \item \textbf{Search:} Google Scholar (scoping), IEEE Xplore (all queries), PubMed and Scopus (targeted) using predefined query strings and result deduplication.
    \item \textbf{Screening:} title/abstract screen, then full‑text eligibility; backward/forward citation chasing on key papers.
    \item \textbf{Extraction:} dataset type (wearable vs clinical), name, availability, task (detection vs prediction), evaluation design (retrospective/pseudo‑prospective/prospective), metrics (sensitivity, FAR/h, AUC), personalization, cohort size/class balance.
    \item \textbf{Quality flags:} sample‑vs alarm‑based evaluation; patient‑preserving splits; cross‑dataset/external validation.
    \item \textbf{Synthesis:} narrative grouping by modality and task with emphasis on ambulatory feasibility and edge constraints.
\end{enumerate}

\section{Key Research Questions}
\begin{enumerate}
    \item Which wearable biosignals (ECG/HRV, EDA, ACC, PPG) yield the best detection and prediction performance across seizure types and settings?
    \item Which model classes (feature‑based vs deep multimodal) generalize across datasets and maintain performance under pseudo‑/prospective, alarm‑based evaluation?
    \item What is the performance gain from patient‑specific fine‑tuning, and how can adaptation be implemented under edge compute and energy constraints?
    \item How can false alarms be minimized while preserving high sensitivity in ambulatory use?
    \item What dataset gaps exist for non‑EEG wearables, and how can augmentation and synthetic data mitigate imbalance without inducing bias?
\end{enumerate}

\section{Overview of Literature Review}
This review is organized as follows.
\begin{enumerate}
    \item \textbf{Modalities:} ECG/HRV, EDA, ACC, PPG; brief physiology, sensor placement, and artifact handling relevant to wearables.
    \item \textbf{Tasks:} seizure detection vs prediction; windowing and labeling conventions; alarm‑based metrics (sensitivity, FAR/h, AUC‑ROC/PR).
    \item \textbf{Models:} classical feature‑based pipelines vs deep 1D‑CNN/LSTM; attention and interpretability considerations.
    \item \textbf{Multimodal fusion:} feature‑, decision‑, and representation‑level fusion strategies and benefits.
    \item \textbf{Personalization:} patient‑specific calibration/fine‑tuning, online adaptation, drift management, and thresholding.
    \item \textbf{Datasets and evaluation:} key datasets and availability; patient‑preserving splits; cross‑dataset transfer; retrospective vs pseudo‑/prospective designs.
    \item \textbf{Edge and deployment:} latency/energy constraints, on‑device inference strategies, robustness in ambulatory settings.
    \item \textbf{Gaps and future directions:} dataset scarcity (non‑EEG wearables), focal seizure detection, prospective validation, standardization, and clinical integration.
\end{enumerate}